% ==========================================
% ТЕМАТИЧЕСКИЙ БАНК: ГИДРОСТАТИКА. СИЛА АРХИМЕДА
% ВАРИАНТЫ 1-10 (ЕГЭ-2026, ДЕМИДОВА)
% ==========================================

\begin{taskbasic}[code={\label{task:fluid_v06_n26}}]
\textbf{Задача 26.} К концам невесомого рычага подвесили за невесомые нерастяжимые нити два сплошных груза массами $1{,}28$~кг и $0{,}32$~кг и привели рычаг в состояние равновесия. Затем рычаг с грузами расположили над водой так, что оба груза целиком оказались под водой. Первоначальное расстояние от более тяжёлого груза до точки опоры равно $20$~см. Определите, на сколько необходимо переместить точку опоры для сохранения равновесия, если объёмы обоих грузов одинаковы и равны $200$~\text{см}$^3$. Сделайте рисунок, на котором укажите силы, действующие на рычаг, для двух случаев, а также на погружённые в воду грузы. Обоснуйте применимость законов, используемых для решения задачи.

\kim{26}
\end{taskbasic}
\answer{0,22}

\begin{taskbasic}[code={\label{task:fluid_v07_n06}}]
\textbf{Задача 6.} Батискаф, полностью погружённый в воду, опускается на дно моря. Как изменяются с течением времени давление воды на дно батискафа и действующая на него архимедова сила? Изменением плотности воды при увеличении глубины пренебречь. Считать, что при погружении батискаф не достиг дна моря.

Для каждой величины определите соответствующий характер изменения:
\begin{enumerate}
    \item увеличивается
    \item уменьшается
    \item не изменяется
\end{enumerate}

Запишите в таблицу выбранные цифры для каждой физической величины. Цифры в ответе могут повторяться.

\begin{center}
\begin{tabular}{|c|c|}
\hline
Давление воды на дно батискафа & Архимедова сила \\ \hline
& \\ \hline
\end{tabular}
\end{center}

\kim{6}
\end{taskbasic}
\answer{13}

\begin{taskbasic}[code={\label{task:fluid_v08_n06}}]
\textbf{Задача 6.} Батискаф, полностью погружённый в воду, поднимается со дна моря с постоянной скоростью. Как изменяются с течением времени кинетическая энергия батискафа и действующая на него архимедова сила? Изменением плотности воды при уменьшении глубины пренебречь. Считать, что в погружённом состоянии батискаф не касался дна моря.

Для каждой величины определите соответствующий характер изменения:
\begin{enumerate}
    \item увеличивается
    \item уменьшается
    \item не изменяется
\end{enumerate}

Запишите в таблицу выбранные цифры для каждой физической величины. Цифры в ответе могут повторяться.

\begin{center}
\begin{tabular}{|c|c|}
\hline
Кинетическая энергия & Архимедова сила \\ \hline
& \\ \hline
\end{tabular}
\end{center}

\kim{6}
\end{taskbasic}
\answer{33}

% ==========================================
% ТЕМАТИЧЕСКИЙ БАНК: ДАВЛЕНИЕ. СИЛА АРХИМЕДА
% ВАРИАНТЫ (ДОПОЛНИТЕЛЬНЫЙ БЛОК)
% ==========================================

\begin{taskbasic}[code={\label{task:fluid_iron_cube}}]
\textbf{Задача 10.} Железный кубик плавает в ртути. Ребро кубика равно $10$~см. Определите силу Архимеда, действующую на кубик.

\kim{4}
\end{taskbasic}
\answer{78}

\begin{taskbasic}[code={\label{task:fluid_copper_cube}}]
\textbf{Задача 11.} Медный кубик, подвешенный на нити, полностью погружён в воду и не касается дна сосуда. Ребро кубика равно $3$~см. Определите силу Архимеда, действующую на кубик.

\kim{4}
\end{taskbasic}
\answer{0,27}

\begin{taskbasic}[code={\label{task:fluid_two_liquids_ball}}]
\textbf{Задача 12.} В стакан налита вода, а поверх неё — керосин. Однородный шар плавает, погружённый в обе жидкости. При этом четверть объёма шара находится в воде. Найдите плотность материала шара.
\begin{center}
    \begin{tikzpicture}[>=Stealth, scale=1.2]
    % --- Штриховка жидкости в сосуде ---
    % Уровень жидкости на высоте 1.8 см, ширина сосуда 2.0 см (от -1 до 1)
    \fill[pattern=dots, pattern color=brandSecondary!60] (-1.0,0) rectangle (1.0,1.8);

    % --- Стенки сосуда и дно ---
    \draw[thick, color=brandSecondary] (-1.0,3.5) -- (-1.0,0) -- (1.0,0) -- (1.0,3.5);

    % --- Поршень сверху ---
    \draw[very thick, color=brandSecondary] (-1.0,2.9) -- (1.0,2.9);

    % --- Плавающий шар ---
    % Центр шара в точке (0, 2.1), радиус R = 0.6 см
    % Заливаем весь шар белым цветом, чтобы полностью перекрыть узор воды внутри него
    \fill[color=white] (0, 2.1) circle (0.6);
    
    % Отрисовка контура шара поверх скрытого уровня воды
    \draw[thick, color=brandSecondary] (0, 2.1) circle (0.6);

    % --- Линии поверхности жидкости (только снаружи шара) ---
    \draw[thick, color=brandSecondary] (-1.0,1.8) -- (-0.565,1.8); 
    \draw[thick, color=brandSecondary] (0.565,1.8) -- (1.0,1.8);
\end{tikzpicture}
\end{center}
\kim{22}
\end{taskbasic}
\answer{850}

\begin{taskbasic}[code={\label{task:fluid_iron_ball_hg}}]
\textbf{Задача 13.} В стакан налита ртуть, а поверх неё — вода. Однородный железный шар плавает, погружённый в обе жидкости. Какая часть объёма шара находится в ртути?
\begin{center}
    \begin{tikzpicture}[>=Stealth, scale=1.2]
    % --- Штриховка жидкости в сосуде ---
    % Уровень жидкости на высоте 1.8 см, ширина сосуда 2.0 см (от -1 до 1)
    \fill[pattern=dots, pattern color=brandSecondary!60] (-1.0,0) rectangle (1.0,1.8);

    % --- Стенки сосуда и дно ---
    \draw[thick, color=brandSecondary] (-1.0,3.5) -- (-1.0,0) -- (1.0,0) -- (1.0,3.5);

    % --- Поршень сверху ---
    \draw[very thick, color=brandSecondary] (-1.0,2.9) -- (1.0,2.9);

    % --- Плавающий шар ---
    % Центр шара в точке (0, 2.1), радиус R = 0.6 см
    % Заливаем весь шар белым цветом, чтобы полностью перекрыть узор воды внутри него
    \fill[color=white] (0, 2.1) circle (0.6);
    
    % Отрисовка контура шара поверх скрытого уровня воды
    \draw[thick, color=brandSecondary] (0, 2.1) circle (0.6);

    % --- Линии поверхности жидкости (только снаружи шара) ---
    \draw[thick, color=brandSecondary] (-1.0,1.8) -- (-0.565,1.8); 
    \draw[thick, color=brandSecondary] (0.565,1.8) -- (1.0,1.8);
\end{tikzpicture}
\end{center}
\kim{22}
\end{taskbasic}
\answer{0,5}

\begin{taskbasic}[code={\label{task:fluid_wood_block_select}}]
\textbf{Задача 14.} Брусок толщиной $6$~см и массой $1$~кг плавает в воде так, что уровень воды приходится на середину бруска. Из приведённого ниже списка выберите все верные утверждения:
\begin{center}

\begin{tikzpicture}[>=Stealth, scale=1.2]
    % --- Пол со штриховкой ---
    \draw[thick, color=brandSecondary] (-1.5,0) -- (1.5,0);
    \fill[pattern=north east lines, pattern color=brandSecondary!50] 
        (-1.5,0) -- (1.5,0) -- (1.5,-0.2) -- (-1.5,-0.2) -- cycle;

    % --- Жидкость в сосуде ---
    % Уровень жидкости на высоте 2.0 см
    \fill[color=brandLight] (-1.0,0) rectangle (1.0,2.0);

    % --- Стенки сосуда ---
    \draw[thick, color=brandSecondary] (-1.0,3.5) -- (-1.0,0) -- (1.0,0) -- (1.0,3.5);

    % --- Плавающий брусок ---
    % Размеры бруска: ширина 1.4 см, высота 0.8 см. Погружен на 0.5 см (от 1.5 до 2.3)
    % Заливаем брусок белым, чтобы перекрыть фон жидкости
    \fill[color=white] (-0.7,1.6) rectangle (0.7,2.4);
    \draw[thick, color=brandSecondary] (-0.7,1.6) rectangle (0.7,2.4);

    % --- Линии поверхности жидкости ---
    % Внешние участки (сплошные)
    \draw[thick, color=brandSecondary] (-1.0,2.0) -- (-0.7,2.0);
    \draw[thick, color=brandSecondary] (0.7,2.0) -- (1.0,2.0);
    
    % Пунктирная линия уровня жидкости внутри бруска
    \draw[thick, dashed, color=brandSecondary] (-0.7,2.0) -- (0.7,2.0);
\end{tikzpicture}
\begin{enumerate}
    \item Если воду заменить на керосин, то глубина погружения бруска увеличится.
    \item Если воду заменить на керосин, то сила Архимеда, действующая на брусок, уменьшится.
    \item Если на брусок положить сверху ещё два таких же бруска, то глубина погружения увеличится на $3$~см.
    \item Плотность материала, из которого изготовлен брусок, равна $500$~кг/м$^3$.
    \item Сила Архимеда, действующая на брусок, равна $10$~Н.
\end{enumerate}
    
\end{center}
\kim{5}
\end{taskbasic}
\answer{145}

\begin{taskbasic}[code={\label{task:fluid_ball_bottom_water}}]
\textbf{Задача 15.} Однородный деревянный шар массой $m = 1{,}6$~кг лежит в сосуде с водой, касаясь дна и не касаясь стенок сосуда, так, что половина шара находится в воде. Определите плотность дерева, если шар давит на дно сосуда с силой $P = 6$~Н. Сделайте рисунок с указанием сил, действующих на шар.

\kim{26}
\end{taskbasic}
\answer{800}

\begin{taskbasic}[code={\label{task:fluid_ball_bottom_kerosene}}]
\textbf{Задача 16.} Однородный деревянный шар лежит в сосуде с керосином, касаясь дна и не касаясь стенок сосуда, так, что половина шара находится в керосине. Определите массу шара, если он давит на дно сосуда с силой $F = 2$~Н. Плотность дерева, из которого сделан шар, равна $\rho = 480$~кг/м$^3$. Сделайте рисунок с указанием сил, действующих на шар.

\kim{26}
\end{taskbasic}
\answer{0,5}

\begin{taskbasic}[code={\label{task:fluid_hollow_ball_lake}}]
\textbf{Задача 17.} Полый стальной шар массой $20$~кг лежит на дне озера. Объём шара равен $15$~дм$^3$. Чему равна сила Архимеда, действующая на шар?

\kim{4}
\end{taskbasic}
\answer{150}

\begin{taskbasic}[code={\label{task:fluid_hollow_ball_float}}]
\textbf{Задача 18.} Полый стальной шар массой $10$~кг плавает на поверхности озера. Объём шара равен $15$~дм$^3$. Чему равна сила Архимеда, действующая на шар?

\kim{4}
\end{taskbasic}
\answer{100}

\begin{taskbasic}[code={\label{task:fluid_ball_water_to_kerosene}}]
\textbf{Задача 19.} Деревянный шарик плавает в воде. Как изменятся глубина погружения шарика и масса вытесненной шариком жидкости, если он будет плавать в керосине?

Для каждой величины определите соответствующий характер изменения:
\begin{enumerate}
    \item увеличится
    \item уменьшится
    \item не изменится
\end{enumerate}

Запишите в таблицу выбранные цифры для каждой физической величины. Цифры в ответе могут повторяться.

\begin{center}
\begin{tabular}{|c|c|}
\hline
Глубина погружения шарика & Масса вытесненной жидкости \\ \hline
& \\ \hline
\end{tabular}
\end{center}

\kim{6}
\end{taskbasic}
\answer{13}

\begin{taskbasic}[code={\label{task:fluid_ball_kerosene_to_water}}]
\textbf{Задача 20.} Деревянный шарик плавает в керосине. Как изменятся сила тяжести и сила Архимеда, действующие на шарик, если он будет плавать в воде?

Для каждой величины определите соответствующий характер изменения:
\begin{enumerate}
    \item увеличивается
    \item уменьшится
    \item не изменится
\end{enumerate}

Запишите в таблицу выбранные цифры для каждой физической величины. Цифры в ответе могут повторяться.

\begin{center}
\begin{tabular}{|c|c|}
\hline
Сила тяжести & Сила Архимеда \\ \hline
& \\ \hline
\end{tabular}
\end{center}

\kim{6}
\end{taskbasic}
\answer{33}


